\section{Introducción a diseños factoriales}

Todos los diseños estudiados hasta ahora en este capítulo estudian el efecto de un \textbf{único} factor de tratamiento (posiblemente controlando fuentes de ruido adicionales mediante bloqueo). Sin embargo, muchas preguntas experimentales reales involucran \textbf{dos o más factores de tratamiento simultáneos} cuyo efecto conjunto —no solo individual— interesa evaluar: por ejemplo, el efecto combinado de la temperatura y la presión sobre el rendimiento de una reacción química, o el efecto conjunto del algoritmo y el tamaño del lote de entrenamiento sobre la exactitud de un modelo de aprendizaje automático.

\begin{definicion}[Diseño factorial completo]
	Un \emph{diseño factorial completo} con dos factores $A$ (con $a$ niveles) y $B$ (con $b$ niveles) consiste en observar \textbf{todas} las $a \times b$ combinaciones posibles de niveles de ambos factores, con $n \geq 1$ réplicas independientes por combinación (para un total de $N = abn$ observaciones). La generalización a más de dos factores sigue el mismo principio: se cruzan todos los niveles de todos los factores.
\end{definicion}

\begin{observacion}
	Frente a la alternativa de variar un factor a la vez manteniendo los demás fijos (\emph{one-factor-at-a-time}), el diseño factorial completo es estrictamente más eficiente: con el mismo número de corridas experimentales, permite estimar no solo los efectos principales de cada factor, sino también sus \textbf{interacciones}, algo que el enfoque secuencial no puede detectar.
\end{observacion}

\subsection{Efectos principales e interacción}

\begin{definicion}[Efecto principal e interacción]
	El \emph{efecto principal} de un factor es el cambio promedio en la variable respuesta al pasar de un nivel del factor a otro, promediando sobre todos los niveles del(los) otro(s) factor(es). Existe \emph{interacción} entre dos factores $A$ y $B$ cuando el efecto de $A$ sobre la variable respuesta \textbf{depende del nivel} de $B$ en el que se evalúe (y viceversa); es decir, las curvas de respuesta de $A$ para distintos niveles de $B$ no son paralelas.
\end{definicion}

El modelo estadístico paramétrico para un diseño factorial completo de dos factores es
\begin{align}
	Y_{ijk} = \mu + \alpha_i + \beta_j + (\alpha\beta)_{ij} + \epsilon_{ijk}, \quad i=1,\ldots,a,\ j=1,\ldots,b,\ k=1,\ldots,n,
	\label{eq:modelo_factorial}
\end{align}
donde $\alpha_i$ y $\beta_j$ son los efectos principales de $A$ y $B$ respectivamente, y $(\alpha\beta)_{ij}$ es el \textbf{efecto de interacción} del par de niveles $(i,j)$, sujeto a las restricciones usuales $\sum_i \alpha_i = \sum_j \beta_j = \sum_i (\alpha\beta)_{ij} = \sum_j (\alpha\beta)_{ij} = 0$.

\begin{teorema}[Partición de la suma de cuadrados en un diseño factorial de dos factores]
	La Suma de Cuadrados Total se descompone en cuatro fuentes ortogonales de variación:
	\begin{align}
		\text{SCT} = \text{SC}_A + \text{SC}_B + \text{SC}_{AB} + \text{SCE},
	\end{align}
	con grados de libertad $\nu_A = a-1$, $\nu_B = b-1$, $\nu_{AB} = (a-1)(b-1)$ para la interacción, y $\nu_E = ab(n-1)$ para el error residual, sobre un total de $\nu_T = abn - 1$.
\end{teorema}

\begin{center}
	\begin{tabular}{lcccc}
		\hline
		\textbf{Fuente de Variación} & \textbf{SC} & \textbf{g.l. ($\nu$)} & \textbf{Cuadrado Medio (CM)} & \textbf{Estadístico $F$} \\ \hline
		Factor $A$ & $\text{SC}_A$ & $a-1$ & $\text{CM}_A = \frac{\text{SC}_A}{a-1}$ & $F_A = \frac{\text{CM}_A}{\text{CME}}$ \\ [1ex]
		Factor $B$ & $\text{SC}_B$ & $b-1$ & $\text{CM}_B = \frac{\text{SC}_B}{b-1}$ & $F_B = \frac{\text{CM}_B}{\text{CME}}$ \\ [1ex]
		Interacción $AB$ & $\text{SC}_{AB}$ & $(a-1)(b-1)$ & $\text{CM}_{AB} = \frac{\text{SC}_{AB}}{(a-1)(b-1)}$ & $F_{AB} = \frac{\text{CM}_{AB}}{\text{CME}}$ \\ [1ex]
		Error residual & $\text{SCE}$ & $ab(n-1)$ & $\text{CME} = \frac{\text{SCE}}{ab(n-1)}$ & \\ [1ex] \hline
		\textbf{Total} & $\text{SCT}$ & $abn-1$ & & \\ \hline
	\end{tabular}
\end{center}

\begin{observacion}[Orden de interpretación]
	Por convención, un análisis factorial siempre se interpreta \textbf{de adentro hacia afuera}: se examina primero el estadístico $F_{AB}$ de la interacción. Si resulta significativo, los efectos principales $\alpha_i$ y $\beta_j$ ya no tienen una interpretación aislada simple (el efecto de $A$ cambia según el nivel de $B$), y el análisis debe reportarse en términos de las combinaciones específicas de niveles. Solo cuando la interacción \textbf{no} es significativa tiene sentido interpretar directamente los efectos principales por separado.
\end{observacion}

\begin{ejemplo}
	\label{exmp:8.6.1}
	Un ingeniero de datos evalúa el tiempo de entrenamiento (en minutos) de un modelo de aprendizaje automático bajo dos algoritmos de optimización ($A_1$, $A_2$) y dos tamaños de lote ($B_1 = 32$, $B_2 = 256$), con $n=3$ réplicas por combinación. Las medias muestrales observadas por celda son:
	\begin{center}
		\begin{tabular}{lcc}
			\hline
			 & $B_1$ (lote 32) & $B_2$ (lote 256) \\ \hline
			$A_1$ & $42$ & $40$ \\
			$A_2$ & $38$ & $20$ \\ \hline
		\end{tabular}
	\end{center}
	¿Sugieren estos promedios la presencia de interacción entre el algoritmo y el tamaño de lote?
\end{ejemplo}

\begin{solucion}
	El efecto de cambiar de $B_1$ a $B_2$ es muy distinto según el algoritmo: bajo $A_1$, el tiempo apenas cambia ($42 \to 40$, una reducción de $2$ minutos), mientras que bajo $A_2$ el tiempo se reduce drásticamente ($38 \to 20$, una reducción de $18$ minutos). Como el efecto del factor $B$ (tamaño de lote) depende marcadamente del nivel del factor $A$ (algoritmo), los datos sugieren una \textbf{interacción sustancial} entre ambos factores. Esto significa que no es correcto afirmar de forma general ``el lote grande reduce el tiempo de entrenamiento''; la afirmación correcta y útil es que el lote grande reduce el tiempo \emph{solo bajo el algoritmo $A_2$}. Confirmar formalmente esta sospecha requeriría calcular $F_{AB}$ y compararlo contra el cuantil crítico $F_{\alpha,\,(a-1)(b-1),\,ab(n-1)}$.
\end{solucion}
